\section{The Grammar of the Liturgical Day}

\subsection{What a liturgical day is}

The code defines its own basic unit before it ranks anything. \rn{4}:

\begin{studyframe}
\latin{Dies liturgicus est dies sanctificatus actionibus liturgicis, praesertim Sacrificio
eucharistico et publica Ecclesiae prece, id est Officio divino; et decurrit a media nocte ad
mediam noctem.}
\medskip

\textit{Working gloss.} A liturgical day is a day sanctified by liturgical actions, especially
by the eucharistic Sacrifice and the public prayer of the Church, that is, the divine Office;
and it runs from midnight to midnight.
\end{studyframe}

Two things follow immediately. The civil day is the container: there is no liturgical day
beginning at sunset in this system. But the \emph{celebration} of a liturgical day is a
different span. \rn{5}: the celebration of a liturgical day runs of itself from Matins to
Compline; there are, however, more solemn days whose Office begins at First Vespers on the
preceding day. And \rn{5} adds a third possibility that will govern the whole treatment of
commemorations: \latin{Habetur denique celebratio liturgica non plena seu sola commemoratio in
Officio et Missa diei liturgici currentis} — there is finally a liturgical celebration that is
not full, or a mere commemoration, in the Office and Mass of the current liturgical day.

The distinction between the day, its celebration, and its First Vespers is what makes two
different collision rules necessary. \rn{6} names what can occupy a day at all: on each day
there is either the Sunday, or the feria, or the vigil, or the feast, or the octave, according
to the calendar and the precedence of liturgical days. Collision between two of these on the
same date is \emph{occurrence}; collision between the Vespers of one day and the First Vespers
of the next is \emph{concurrence}. Both are treated in Section~\ref{sec:precedence}.

\subsection{Four classes, and what they replaced}

\rn{8} is one sentence: \latin{Dies liturgici sunt primae, secundae, tertiae aut quartae
classis} — liturgical days are of the first, second, third, or fourth class. This is the single
most visible change of 1960, and it is worth being exact about what it did and did not do.

It did not invent four ranks out of nothing. It abolished a grading vocabulary — \latin{duplex
I classis}, \latin{duplex II classis}, \latin{duplex maius}, \latin{duplex minus},
\latin{semiduplex}, \latin{simplex} — that had come to carry two jobs at once: a rank in the
order of precedence, and a description of how the Office was structured (whether antiphons were
doubled, how many lessons at Matins, whether the day had First Vespers). By 1960 the second job
had largely been separated out, and the words had become opaque labels. The code replaced the
labels with ordinals and located the entire precedence question in one table.

The conversion is stated arithmetically in the \latin{Variationes in Calendario}, \rnn{1--4},
and it is a mapping, not a re-evaluation of individual feasts:

\begin{center}
\small
\begin{tabular}{@{}>{\raggedright\arraybackslash}p{0.42\linewidth}>{\raggedright\arraybackslash}p{0.30\linewidth}@{}}
\toprule
\textbf{Old grade as inscribed in calendars} & \textbf{New grade from 1961} \\
\midrule
\latin{duplex I classis} & feast of the I class \\
\latin{duplex II classis} & feast of the II class \\
\latin{duplex maius} or \latin{duplex minus} & feast of the III class \\
\latin{semiduplex} (from 1955, \latin{simplex}) & feast of the III class \\
\latin{simplex}, already reduced to a commemoration in 1955 & inscribed as a commemoration \\
\bottomrule
\end{tabular}
\end{center}

The fourth line repays attention. The general decree of the Sacred Congregation of Rites of
23 March 1955 — which \latin{Rubricarum instructum} \rn{2} names only by its date and its title
\latin{De rubricis ad simpliciorem formam redigendis} — had already suppressed the semidouble
grade by reducing semidoubles to simples; the 1960 code then promoted that whole stratum to the
III class. The fifth line is the other half of the same operation:
feasts that were simple before 1955 and had been reduced to commemorations then are now
inscribed in the calendar as commemorations, which under the new law is not a rank of feast at
all but a mode of non-full celebration. Ten further feasts were reduced to commemorations by
\rn{5} of the \latin{Variationes} itself; they are listed in Section~\ref{sec:sanctoral}.

\subsection{The classes are not uniform across kinds of day}

A common error is to read ``II class'' as a single grade. It is not. The code assigns classes
separately to Sundays, ferias, vigils, and feasts, and the ranges differ:

\begin{center}
\small
\begin{tabular}{@{}l l l@{}}
\toprule
\textbf{Kind of day} & \textbf{Classes available} & \textbf{Locus} \\
\midrule
Sunday & I or II & \rn{10} \\
Feria & I, II, III, or IV & \rn{22} \\
Vigil & I, II, or III & \rn{29} \\
Feast & I, II, or III & \rn{36} \\
Octave & I or II & \rn{65} \\
\bottomrule
\end{tabular}
\end{center}

There is no Sunday of the III class, no feast of the IV class, and no vigil of the IV class.
The IV class exists only for ferias and, through \rn{78}, for the Office of Holy Mary on
Saturday. Conversely, the class alone never settles a conflict between two days of unequal
kind: a II class Sunday and a II class feast are both ``II class'', and the table at \rn{91}
places the Sunday above ordinary universal II class feasts but below II class feasts of the
Lord. This is exactly what \rn{7} means by saying that precedence is determined solely by the
table.

\subsection{First Vespers, and who has them}

\rn{37} states how feasts are celebrated, and in doing so fixes which days reach back into the
preceding evening:

\begin{studybox}{\rn{37}: the manner of celebrating feasts}
\begin{enumerate}[label=\alph*), leftmargin=1.6em, itemsep=0.2em]
\item feasts of the I class are counted among the more solemn days, whose Office begins at
First Vespers on the preceding day;
\item feasts of the II and III class have an Office that runs of itself from Matins to Compline
of the day itself;
\item feasts of the Lord of the II class \emph{acquire} First Vespers whenever, in occurrence,
they take the place of a Sunday of the II class.
\end{enumerate}
\end{studybox}

Clause (c) is the hinge on which many concurrence questions turn, and it is conditional: a II
class feast of the Lord has no First Vespers of its own, but gains them by taking a Sunday's
place under \rn{16}\,(a). Sundays always have First Vespers by \rn{13}; the Office of a Sunday
begins at First Vespers on the preceding Saturday and ends after Compline of the Sunday.

\subsection{Universal, particular, proper, indult}

\rn{38} divides feasts on a second axis that is independent of class: \latin{Festa sunt
universalia aut particularia; particularia autem sunt propria aut indulta}. Universal feasts
are those the Holy See inscribes in the calendar of the universal Church (\rn{39}) and must be
celebrated by all who follow the Roman Rite. Particular feasts are inscribed in particular
calendars either by law or by indult of the Holy See (\rn{40}), must be celebrated by all bound
to that calendar, and — a provision often forgotten — \latin{nonnisi ex speciali Sanctae Sedis
indulto e calendario expungi vel gradu mutari possunt}: they can be struck from the calendar or
changed in grade only by special indult of the Holy See.

\rnn{41--46} then enumerate exhaustively which proper feasts enter a particular calendar by law
and at what class. The list is short enough to give whole, because it is the entire legal
content of the phrase ``proper calendar'' in this system:

\begingroup\small
\renewcommand{\arraystretch}{1.12}
\begin{longtable}{@{}>{\raggedright\arraybackslash}p{0.28\linewidth}>{\raggedright\arraybackslash}p{0.46\linewidth}>{\raggedright\arraybackslash}p{0.14\linewidth}@{}}
\toprule
\textbf{Owner} & \textbf{Proper feast} & \textbf{Class} \\
\midrule
\endfirsthead
\toprule
\textbf{Owner} & \textbf{Proper feast} & \textbf{Class} \\
\midrule
\endhead
\bottomrule
\endfoot
Nation, region or province, whether ecclesiastical or civil (\rn{42}) & principal Patron duly constituted & I \\*
 & secondary Patron duly constituted & II \\
\midrule
Diocese or other territory under a local Ordinary (\rn{43}) & principal Patron duly constituted & I \\*
 & anniversary of the dedication of the cathedral & I \\*
 & secondary Patron duly constituted & II \\*
 & Saints and Blessed duly inscribed in the Martyrology or its Appendix having a particular relation to the diocese, such as origin, longer residence, or death & II or III, or a commemoration \\
\midrule
Place, town, or city (\rn{44}) & principal Patron duly constituted & I \\*
 & secondary Patron duly constituted & II \\
\midrule
Church, or public or semi-public oratory taking the place of a church (\rn{45}) & anniversary of dedication, if consecrated & I \\*
 & feast of the Title, if consecrated or at least solemnly blessed & I \\*
 & feast of a Saint whose body is kept there & II \\*
 & feast of a Blessed whose body is kept there & III \\
\midrule
Order or Congregation (\rn{46}) & feast of the Title & I \\*
 & canonised Founder & I \\*
 & beatified Founder & II \\*
 & principal Patron of the whole Order or Congregation, and of a religious province in that province & I \\*
 & secondary Patron, likewise & II \\*
 & Saints and Blessed who were members & II or III, or a commemoration \\
\end{longtable}
\endgroup


Everything else in a particular calendar is an \latin{indultum} — a feast inscribed by grant of
the Holy See, on the day the grant assigns (\rnn{47}, 61\,d). The distinction is not decorative:
in the table of precedence, proper feasts and indult feasts of the same class occupy different
lines, and among indult feasts movable ones precede fixed ones (lines 13, 20, 23).

\subsection{The calendars themselves}

\rnn{48--58} build the calendar system in layers. The universal calendar is the one prefixed to
the Roman Breviary and Missal (\rn{49}). A particular calendar — diocesan or religious — is made
\latin{inserendo calendario universali festa particularia}, by inserting particular feasts into
the universal calendar (\rn{50}); it is drawn up by the local Ordinary or by the supreme
Moderator of the religion with the counsel of his Chapter or general Council, and \emph{must be
approved} by the Sacred Congregation of Rites. On the diocesan calendar are then built the
calendar of each place, of each church or oratory, and of certain religious institutes
(\rn{53}); on the religious calendar, those of each religious province and of each religious
house (\rn{56}).

Two obligations cut across the layers and are easy to miss. By \rn{57}, religious in every
diocese and place — \latin{ii etiam qui alium ritum ac romanum sequuntur}, even those following
a rite other than the Roman — must celebrate together with the diocesan clergy the principal
Patron of the nation, region or province, diocese, place, town, or city (I class), the
anniversary of the dedication of the cathedral (I class), and any other feasts that are actually
days of precept, at the grade at which the diocesan calendar inscribes them. By \rn{58},
religious must conform to the diocesan clergy as to day and more proper Office where their own
Saints are honoured as principal Patrons, and may adopt a higher grade or more proper Office
where the local clergy celebrate their Saints so, provided the feasts stand on the same day in
both calendars.

The Missal rubrics complete the picture from the celebrant's side. In a church or public
oratory every priest, diocesan or religious, must celebrate according to the calendar of that
church or oratory (Missal rubrics, \rn{275}); in secondary oratories he may follow either that
oratory's calendar or his own (\rn{276}); in private oratories and on a portable altar outside
a sacred place, either the calendar of the place or his own (\rn{277}); an oratory fixed in a
ship is a public oratory and uses the calendar of the universal Church (\rn{279}). And by
\rn{278}, whatever calendar he might otherwise follow, every priest must say the Mass of the
principal Patron of the nation, region or province, diocese, town, or city, of the anniversary
of the dedication of the cathedral, and of other feasts that are actually days of precept.
